\clearpage
\cleardoublepage

\chapter{Classic Network Architectures}

\section{Evolution of Neural Networks}

The history of deep learning shows a progression from simple networks to increasingly complex architectures, each building upon innovations of the previous era.

\section{Perceptron and Early Neural Networks}

\subsection{The Perceptron (1958)}

The perceptron, introduced by Frank Rosenblatt, was the first attempt at artificial neural networks:
\begin{equation}
y = \begin{cases} 1 & \text{if } \mathbf{w} \cdot \mathbf{x} + b > 0 \\ 0 & \text{otherwise} \end{cases}
\end{equation}

Limitations: Can only solve linearly separable problems (XOR problem showed this).

\section{LeNet (1998)}

Yann LeCun's LeNet was the first successful CNN for digit recognition.

\begin{architecture}[LeNet-5 Architecture]
\begin{enumerate}
    \item Input: $32\times32$ grayscale image
    \item C1: $6$ filters of $5\times5$, stride 1 $\rightarrow$ $28\times28\times6$
    \item S2: $2\times2$ avg pool, stride 2 $\rightarrow$ $14\times14\times6$
    \item C3: $16$ filters of $5\times5$ $\rightarrow$ $10\times10\times16$
    \item S4: $2\times2$ avg pool $\rightarrow$ $5\times5\times16$
    \item C5: $120$ filters of $5\times5$ $\rightarrow$ $120$
    \item F6: Fully connected, 84 units
    \item Output: 10 classes (digits)
\end{enumerate}
\end{architecture}

Total parameters: $\sim$60,000

Key innovations:
\begin{itemize}
    \item Weight sharing through convolutions
    \item Local receptive fields
    \item Spatial subsampling (pooling)
\end{itemize}

\section{AlexNet (2012)}

AlexNet won ImageNet LSVRC-2012 with a significant margin, revolutionizing computer vision.

\begin{architecture}[AlexNet Architecture]
\begin{enumerate}
    \item Input: $224\times224\times3$ RGB image
    \item Conv1: $96$ filters of $11\times11$, stride 4 $\rightarrow$ $55\times55\times96$
    \item Pool1: $3\times3$ max pool, stride 2 $\rightarrow$ $27\times27\times96$
    \item Norm1: Response normalization
    \item Conv2: $256$ filters of $5\times5$ $\rightarrow$ $27\times27\times256$
    \item Pool2: $3\times3$ max pool $\rightarrow$ $13\times13\times256$
    \item Conv3: $384$ filters of $3\times3$
    \item Conv4: $384$ filters of $3\times3$
    \item Conv5: $256$ filters of $3\times3$
    \item Pool5: $3\times3$ max pool $\rightarrow$ $6\times6\times256$
    \item FC6: 4096 neurons
    \item FC7: 4096 neurons
    \item FC8: 1000 classes (ImageNet)
\end{enumerate}
\end{architecture}

Key innovations:
\begin{itemize}
    \item ReLU activation (faster training)
    \item GPU training (5-6 days vs weeks)
    \item Data augmentation
    \item Dropout regularization
    \item Response normalization
\end{itemize}

Parameters: $\sim$60 million

\section{VGGNet (2014)}

VGGNet emphasized depth and simplicity, using only $3\times3$ convolutions.

\begin{architecture}[VGG-16 Architecture]
\begin{enumerate}
    \item Input: $224\times224\times3$
    \item Conv blocks:
    \begin{itemize}
        \item Block 1: $2\times$ conv $64$, pool
        \item Block 2: $2\times$ conv $128$, pool
        \item Block 3: $3\times$ conv $256$, pool
        \item Block 4: $3\times$ conv $512$, pool
        \item Block 5: $3\times$ conv $512$, pool
    \end{itemize}
    \item FC: 4096, 4096, 1000
\end{enumerate}
\end{architecture}

Key insight: Two $3\times3$ convolutions have receptive field of $5\times5$ but fewer parameters:
\begin{itemize}
    \item $5\times5$ conv: $25C^2$ parameters
    \item Two $3\times3$ convs: $18C^2$ parameters (28\% reduction)
    \item More nonlinearities with two ReLUs
\end{itemize}

Parameters: $\sim$138 million

\section{GoogLeNet / Inception (2014)}

Introduced the Inception module with parallel multi-scale convolutions.

\begin{architecture}[Inception Module]
Parallel operations at different scales:
\begin{itemize}
    \item $1\times1$ conv (bottleneck)
    \item $3\times3$ conv (after $1\times1$)
    \item $5\times5$ conv (after $1\times1$)
    \item $3\times3$ max pool (after $1\times1$)
\end{itemize}
Concatenate all outputs along channel dimension.
\end{architecture}

Key innovations:
\begin{itemize}
    \item Multi-scale feature extraction
    \item Dimensionality reduction with $1\times1$ conv
    \item Global Average Pooling
    \item Auxiliary classifiers for training
\end{itemize}

Parameters: $\sim$4 million (very efficient!)

\section{ResNet (2015)}

ResNet introduced residual connections, enabling training of very deep networks.

\begin{architecture}[ResNet-50 (Bottleneck)]
\begin{enumerate}
    \item Conv1: $7\times7$, 64, stride 2
    \item Pool: $3\times3$ max, stride 2
    \item Stage 1: 3 blocks, 64 channels
    \item Stage 2: 4 blocks, 128 channels
    \item Stage 3: 6 blocks, 256 channels
    \item Stage 4: 3 blocks, 512 channels
    \item Pool: Global Average
    \item FC: 1000 classes
\end{enumerate}
\end{architecture}

Key innovations:
\begin{itemize}
    \item Skip connections for gradient flow
    \item Identity shortcuts
    \item Bottleneck design ($1\times1$, $3\times3$, $1\times1$)
    \item Enables 1000+ layers
\end{itemize}

Parameters: $\sim$25 million

\section{DenseNet (2017)}

Dense connections where each layer receives inputs from all preceding layers.

\begin{architecture}[DenseNet Block]
For a block with $L$ layers, each receives $k_0 + k \times (l-1)$ inputs:
\begin{align}
\mathbf{x}_1 &= H_1([\mathbf{x}_0]) \\
\mathbf{x}_2 &= H_2([\mathbf{x}_0, \mathbf{x}_1]) \\
\mathbf{x}_3 &= H_3([\mathbf{x}_0, \mathbf{x}_1, \mathbf{x}_2]) \\
\vdots \\
\mathbf{x}_L &= H_L([\mathbf{x}_0, \mathbf{x}_1, \ldots, \mathbf{x}_{L-1}])
\end{align}
\end{architecture}

Key innovations:
\begin{itemize}
    \item Feature reuse reduces parameters
    \item Strong gradient flow
    \item Better parameter efficiency than ResNet
\end{itemize}

Parameters: $\sim$27 million (similar to ResNet but better)

\section{Transformer (2017)}

The Transformer architecture replaced recurrence with self-attention.

\begin{architecture}[Transformer Architecture]
\begin{enumerate}
    \item \textbf{Encoder:} $N$ identical layers
    \begin{itemize}
        \item Multi-head self-attention
        \item Feed-forward network
        \item Layer normalization
        \item Residual connections
    \end{itemize}
    \item \textbf{Decoder:} $N$ identical layers
    \begin{itemize}
        \item Masked multi-head self-attention
        \item Encoder-decoder attention
        \item Feed-forward network
    \end{itemize}
    \item \textbf{Positional encoding} (sinusoidal or learned)
\end{enumerate}
\end{architecture}

Key innovations:
\begin{itemize}
    \item Self-attention replaces recurrence
    \item Parallelizable (unlike RNNs)
    \item Captures long-range dependencies
    \item Scalable with computation
\end{itemize}

Parameters: 65M (base) to 175B (GPT-3) and beyond

\section{MobileNet (2017)}

MobileNet introduced depthwise separable convolutions for efficient mobile inference (Howard et al., 2017).

\begin{architecture}[MobileNet Building Block]
\begin{enumerate}
    \item Depthwise convolution: $K \times K$ conv, 1 filter per input channel
    \item Pointwise convolution: $1 \times 1$ conv, $C_{out}$ filters
\end{enumerate}
FLOPs reduction vs standard convolution:
\begin{equation}
\frac{\text{Standard}}{\text{Separable}} = \frac{K^2 \cdot C_{in} \cdot C_{out}}{K^2 \cdot C_{in} + C_{in} \cdot C_{out}} \approx \frac{1}{\frac{1}{C_{out}} + \frac{1}{K^2}}
\end{equation}
For $K = 3$, $C_{out} = 256$: speedup $\approx 8$--$9\times$.
\end{architecture}

\begin{properties}[MobileNet Design Choices]
\begin{itemize}
    \item \textbf{Width multiplier} $\alpha$: Reduces channels uniformly ($\alpha \in \{0.25, 0.5, 0.75, 1.0\}$)
    \item \textbf{Resolution multiplier} $\rho$: Reduces input resolution ($224, 192, 160, 128$)
    \item MobileNetV1: 4.2M params, 0.575 GMACs for $\alpha = 1.0$
    \item MobileNetV2: Adds inverted residuals and linear bottlenecks (3.4M params)
    \item MobileNetV3: NAS-optimized, hardware-aware design (2.5M--5.4M params)
\end{itemize}
\end{properties}

\section{SENet (2018)}

Squeeze-and-Excitation Networks (Hu et al., 2018) introduce channel attention to adaptively recalibrate feature maps.

\begin{architecture}[SE Block]
\begin{enumerate}
    \item \textbf{Squeeze}: Global average pooling $\rightarrow$ $1 \times 1 \times C$ descriptor
    \item \textbf{Excitation}: Two FC layers with bottleneck (reduction ratio $r = 16$)
    \begin{align}
    \mathbf{z} &= \mathbf{W}_2 \cdot \delta(\mathbf{W}_1 \cdot \mathbf{s}) \\
    \text{where}\quad \mathbf{s} &= \frac{1}{H \cdot W}\sum_{i,j} \mathbf{X}_{i,j} \quad \text{(squeeze)}
    \end{align}
    \item \textbf{Scale}: Channel-wise multiplication $\tilde{\mathbf{X}}_c = \mathbf{X}_c \cdot \sigma(\mathbf{z})_c$
\end{enumerate}
\end{architecture}

Key insight: Different channels have different importance for different inputs. The SE block learns which channels to emphasize or suppress. Added overhead is minimal ($< 1\%$ of total FLOPs) but consistently improves accuracy by 1--2\%.

\section{EfficientNet (2019)}

EfficientNet (Tan \& Le, 2019) scales all three dimensions (depth, width, resolution) simultaneously using a compound scaling coefficient $\phi$:
\begin{align}
\text{depth:} \quad d &= \alpha^\phi \\
\text{width:} \quad w &= \beta^\phi \\
\text{resolution:} \quad r &= \gamma^\phi \\
\text{subject to:} \quad \alpha \cdot \beta^2 \cdot \gamma^2 &\approx 2 \quad \text{(target FLOPs $\approx 2^\phi$)}
\end{align}

\begin{table}[ht]
    \centering
    \caption{EfficientNet Family}
    \begin{tabular}{L{2cm} L{2cm} L{2cm} L{3cm} L{2.5cm}}
        \toprule
        \textbf{Variant} & \textbf{Params} & \textbf{FLOPs} & \textbf{Resolution} & \textbf{Top-1 Acc.} \\
        \midrule
        B0 & 5.3M & 0.39B & $224^2$ & 77.1\% \\
        B1 & 7.8M & 0.70B & $240^2$ & 79.1\% \\
        B2 & 9.2M & 1.0B & $260^2$ & 80.1\% \\
        B3 & 12M & 1.8B & $300^2$ & 81.6\% \\
        B4 & 19M & 4.2B & $380^2$ & 82.9\% \\
        B7 & 66M & 37B & $600^2$ & 84.3\% \\
        \bottomrule
    \end{tabular}
\end{table}

EfficientNet achieves state-of-the-art accuracy with significantly fewer parameters and FLOPs than previous architectures. The compound scaling principle has become a standard design guideline.

\section{Vision Transformer (ViT) (2020)}

Vision Transformer (Dosovitskiy et al., 2020) applies the transformer architecture directly to images, treating patches as tokens.

\begin{architecture}[ViT Architecture]
\begin{enumerate}
    \item \textbf{Patch embedding}: Split $224 \times 224$ image into $16 \times 16$ patches ($196$ tokens), linearly project each
    \item \textbf{[CLS] token}: Prepend a learnable class token
    \item \textbf{Positional encoding}: Add learned position embeddings
    \item \textbf{Transformer encoder}: $L$ layers of multi-head self-attention + MLP
    \item \textbf{Classification head}: Use [CLS] token output for classification
\end{enumerate}
\end{architecture}

\begin{properties}[ViT Key Insights]
\begin{itemize}
    \item \textbf{No inductive bias}: Unlike CNNs, ViT has no built-in translation equivariance or locality --- it must learn these from data
    \item \textbf{Data hungry}: ViT requires large-scale pretraining (ImageNet-21k or JFT-300M) to outperform CNNs
    \item \textbf{Scales well}: Performance improves consistently with more data and compute
    \item \textbf{ViT-B/16}: 86M params, 17.6B FLOPs, 79.9\% Top-1 (ImageNet, from 21k pretrain)
    \item \textbf{ViT-L/16}: 307M params, 61.6B FLOPs, 82.3\% Top-1
    \item \textbf{DeiT}: Distillation enables ViT training from scratch on ImageNet-1k alone
\end{itemize}
\end{properties}

\begin{note}[CNNs vs ViTs: The Current Landscape]
As of 2024, the field has converged toward hybrid approaches:
\begin{itemize}
    \item \textbf{ConvNeXt}: Modern CNNs redesigned to match ViT scalability (Liu et al., 2022)
    \item \textbf{Swin Transformer}: Hierarchical ViT with shifted windows for dense prediction
    \item \textbf{Both paradigms achieve similar performance}; the choice often depends on downstream task and infrastructure
\end{itemize}
\end{note}

\section{U-Net (2015)}

U-Net (Ronneberger et al., 2015) is an encoder-decoder architecture designed for biomedical image segmentation. Its U-shaped structure gives it its name.

\begin{architecture}[U-Net Architecture]
\begin{enumerate}
    \item \textbf{Contracting path (encoder)}: 4 blocks of $3\times3$ conv + ReLU + $2\times2$ max pool
    \item \textbf{Bottleneck}: $3\times3$ conv at the lowest resolution
    \item \textbf{Expansive path (decoder)}: Upsampling + concatenation with skip connections + $3\times3$ conv
    \item Skip connections connect each encoder level to the corresponding decoder level
\end{enumerate}
\end{architecture}

\begin{properties}[U-Net Design Principles]
\begin{itemize}
    \item \textbf{Skip connections}: Preserve fine-grained spatial detail lost during downsampling
    \item \textbf{Symmetric design}: Encoder and decoder have matching resolutions
    \item \textbf{Works with limited data}: Original paper achieved excellent results with only 30 training images
    \item Widely extended: Attention U-Net, U-Net++, nnU-Net (automated configuration)
    \item Foundation of modern segmentation: used in medical imaging, satellite analysis, and autonomous driving
\end{itemize}
\end{properties}

\section{ShuffleNet (2018)}

ShuffleNet (Zhang et al., 2018) introduces channel shuffling to enable efficient pointwise group convolutions.

The key operation, \textbf{channel shuffle}, rearranges channels across groups before the next group convolution:
\begin{equation}
\text{Shuffle}(\mathbf{X})_{g, c, h, w} = \mathbf{X}_{\lfloor c / (C/G) \rfloor, \bigl(g + c \cdot G \bigr) \bmod C, h, w}
\end{equation}

Without shuffling, stacked group convolutions would only communicate within their group (blocking cross-group information flow). Channel shuffle provides an efficient alternative to the $1 \times 1$ convolution used in MobileNet.

\begin{properties}[ShuffleNet vs MobileNet]
\begin{itemize}
    \item ShuffleNetV1: 1.5$\times$ faster than MobileNet on ARM devices
    \item ShuffleNetV2: Introduces channel split and four guidelines for efficient design (equal channel width, balanced ops, minimize memory access, element-wise ops cheap)
    \item ShuffleNetV2 ($1.5\times$): 3.5M params, 299M MACs, 73.3\% Top-1
\end{itemize}
\end{properties}

\begin{table}[ht]
    \centering
    \caption{Architecture Comparison}
    \begin{tabular}{L{2.5cm} L{1.5cm} L{2cm} L{2.5cm} L{3cm}}
        \toprule
        \textbf{Network} & \textbf{Year} & \textbf{Parameters} & \textbf{Top-1 Acc.} & \textbf{Key Innovation} \\
        \midrule
        AlexNet & 2012 & 60M & 62.5\% & GPU training, ReLU \\
        VGG-16 & 2014 & 138M & 71.5\% & $3\times3$ conv stack \\
        GoogLeNet & 2014 & 4M & 72.0\% & Inception module \\
        ResNet-50 & 2015 & 25M & 76.1\% & Skip connections \\
        DenseNet-121 & 2017 & 8M & 74.4\% & Dense connections \\
        MobileNetV2 & 2018 & 3.4M & 72.0\% & Depthwise separable \\
        EfficientNet-B7 & 2019 & 66M & 84.3\% & Compound scaling \\
        ViT-B/16 & 2021 & 86M & 79.9\% & Patch-based attention \\
        ConvNeXt-L & 2022 & 350M & 85.1\% & Modernized CNN \\
        \bottomrule
    \end{tabular}
\end{table}

\section{Chapter Summary}

\begin{enumerate}
    \item LeNet pioneered CNNs for digit recognition
    \item AlexNet popularized deep learning with GPU training and ReLU
    \item VGG showed the power of depth through $3\times3$ convolution stacks
    \item Inception introduced multi-scale processing and dimensionality reduction
    \item ResNet enabled ultra-deep networks with skip connections
    \item DenseNet improved parameter efficiency through feature reuse
    \item MobileNet and ShuffleNet brought efficient inference to mobile devices
    \item SENet introduced channel attention, later extended to CBAM
    \item EfficientNet showed that compound scaling outperforms single-dimension scaling
    \item ViT demonstrated that pure attention can match or exceed CNNs at scale
    \item U-Net defined the encoder-decoder paradigm for segmentation
    \item The field continues to converge: modern CNNs and ViTs achieve similar performance
\end{enumerate}

\section{Further Reading}

\begin{itemize}
    \item Krizhevsky, A., Sutskever, I., \& Hinton, G. (2012). ImageNet classification with deep convolutional neural networks. \textit{NeurIPS}.
    \item He, K., Zhang, X., Ren, S., \& Sun, J. (2016). Deep residual learning for image recognition. \textit{CVPR}.
    \item Hu, J., Shen, L., \& Sun, G. (2018). Squeeze-and-excitation networks. \textit{CVPR}.
    \item Howard, A. et al. (2017). MobileNets: Efficient convolutional neural networks for mobile vision applications. \textit{arXiv:1704.04861}.
    \item Tan, M. \& Le, Q. (2019). EfficientNet: Rethinking model scaling for convolutional neural networks. \textit{ICML}.
    \item Dosovitskiy, A. et al. (2021). An image is worth 16x16 words: Transformers for image recognition at scale. \textit{ICLR}.
    \item Ronneberger, O., Fischer, P., \& Brox, T. (2015). U-Net: Convolutional networks for biomedical image segmentation. \textit{MICCAI}.
    \item Liu, Z. et al. (2022). A convnet for the 2020s. \textit{CVPR}.
\end{itemize}
